Any unit propagation process is performed following the \textbf{watched literals} behavior. Watched literals is the best known way to approximate the 
naive algorithm. The approximated algorithm is defined by the following steps:
\begin{itemize}
    \renewcommand{\labelitemi}{-}
    \item Watch two unassigned literals in every clause in preparation of unit propagation.
    \item Examine the clause only when one of them is assigned to false.
    \item For each literal, maintain a watch list containing all the clauses that are composed by that literal.
\end{itemize}
Whenever a literal $l$ is assigned to true, for each clause $k$ that currently is watching or observing the negated literal $\neg l$, we do:
\begin{itemize}
    \renewcommand{\labelitemi}{-}
    \item If all but one literal $l$ is assigned to false, assign $l$ to true.
    \item If all literals are assigned to false, return the unsatisfiability of the problem.
    \item If any literal is assigned to true, continue.
    \item Otherwise, remove $k$ from the watch list of $\neg l$ and add it to the watch list of one of its remaining unassigned literals.
\end{itemize}
\begin{example}
    e.g. Watched literals mechanism. \vspace{7pt}

    Given the following clause:
    $$ c: x_1 \vee \neg x_2 \vee x_4 \vee x_4 $$
    We are gonna to introduce the main considerations about the watched literals behavior. We suppose that initially the approximated algorithm decided to 
    observe $x_1$ and $\neg x_2$ literals from clause $c$. \vspace{7pt}

    \begin{itemize}
        \renewcommand{\labelitemi}{-} 
        \item $x_3$ is assigned to false ($\neg x_3$), nothing happens, since the clause $c$ is not in the watch list of the boolean variable $x_3$.
        \item $x_2$ is assign to true ($x_2 \rightarrow \neg x_2 \rightarrow \text{false}$), therefore clause $c$ is inspected. Since there is another literal to observe, we remove $c$ from the watch list of $\neg x_2$ and we add the same clause $c$ in the watch list of $x_4$.
    \end{itemize} \vspace{7pt}

    The current situation is:
    $$ c: x_1 \vee \ \text{false} \ \vee \ \text{false} \ \vee x_4 $$
    both last literals are observed by the clause $c$. \vspace{7pt}

    \begin{itemize}
        \renewcommand{\labelitemi}{-} 
        \item $x_4$ is assigned to false ($\neg x_4$). Clause $c$ is again inspected, since it is already in the watch list of $x_4$, so we do unit propagation: all the literals of clause $c$ are assigned to false, so we assign the last literal to true.
        \item $x_1$ is assigned to true, stating the satisfiability of the clause.
    \end{itemize}
\end{example}